Let $p(n)$ denote the least prime factor of $n$ (so if $n$ is prime then $p(n)=n$).
The extracted problem states that there exists a constant $c>0$ such that
\[S(x):=\sum_{\substack{n<x\\ n\ \textrm{not prime}}}\frac{p(n)}{n}\ \sim\ c\,\frac{x^{1/2}}{(\log x)^2}.
\]
It then asks whether there exists a constant $C>0$ such that
\[T(x):=\sum_{x\le n\le x+C x^{1/2}(\log x)^2}\frac{p(n)}{n}\ \gg\ 1\]
for all sufficiently large $x$.

\textbf{Ambiguity.} In the second sum, primes $n$ are \emph{not} explicitly excluded. If primes are included, then each prime contributes $p(n)/n=1$ and the question becomes closely related to guaranteeing primes in intervals of length $\asymp x^{1/2}(\log x)^2$ (a Legendre-type prime gap question).  The reference is ambiguous about whether primes should be excluded.  In what follows I focus on the \emph{corrected} version consistent with the first sum:
\[T_{\mathrm{comp}}(x):=\sum_{\substack{x\le n\le x+C x^{1/2}(\log x)^2\\ n\ \textrm{not prime}}}\frac{p(n)}{n},\]
and ask whether $T_{\mathrm{comp}}(x)\gg 1$ for all large $x$.
